{\scriptsize\setlength{\tabcolsep}{3pt}
\begin{longtable}{rcccccp{5.7cm}}
\caption{Casos de teste de \texttt{stack.vwf}: visão geral.}
\label{tab:caso-stack}\\
\hline $t$ (ns) & \texttt{nrst} & \texttt{stack\_in} & \texttt{stack\_push} & \texttt{stack\_pop} & \texttt{stack\_out} & Observação \\ \hline \endfirsthead
\hline $t$ (ns) & \texttt{nrst} & \texttt{stack\_in} & \texttt{stack\_push} & \texttt{stack\_pop} & \texttt{stack\_out} & Observação \\ \hline \endhead
0--20 & 0 & 0 & 0 & 0 & 0 &  \\
20--40 & 1 & 100 & 1 & 0 & 100 & solta o reset; push 100 \\
40--60 & 1 & 200 & 1 & 0 & 200 & push 200 \\
60--80 & 1 & 300 & 1 & 0 & 300 & push 300 \\
80--100 & 1 & 999 & 0 & 1 & 200 & pop (volta a 200) \\
100--120 & 1 & 999 & 0 & 1 & 100 & pop (volta a 100) \\
120--140 & 1 & 400 & 1 & 0 & 400 & push 400 \\
140--160 & 1 & 555 & 1 & 1 & 555 & push e pop juntos: push 555 vence \\
160--180 & 1 & 777 & 0 & 0 & 555 & sem operação (stack\_in muda e é ignorado) \\
180--200 & 1 & 777 & 0 & 1 & 400 & pop (555 sai, volta a 400) \\
200--220 & 1 & 777 & 0 & 1 & 100 & pop (volta a 100) \\
220--240 & 1 & 777 & 0 & 1 & 0 & pop (esvazia, topo = 0) \\
240--260 & 1 & 777 & 0 & 1 & 0 & pop com a pilha vazia: continua 0 \\
260--280 & 1 & 777 & 0 & 0 & 0 & sem operação \\
\hline
\end{longtable}}
\vspace{-0.9cm}
\fonte{Autoria própria. Saídas do fim de cada intervalo, conforme o modelo de referência e a simulação do Quartus II.}
